\subsection{Simulation Interrupt\slash Control Files}
\label{SEC:Files:Interrupt}

There are several temporary files that can be used to interrupt a running
simulation to request a certain action or modify certain parameters while the
simulation continues to run.
%
This is particularly useful on large HPC systems where the alternative ---
stopping the simulation, editing the input file, then resubmitting the job to
the system queue, and waiting for it to start --- could take days or even weeks
for large jobs needing thousands of MPI processes.
%
After completion of every $N\supth$ time step, GFR checks the current run
directory for the presence of a file with one of the following case-sensitive
file names:
%
\begin{itemize}
  \item \nameref{SEC:Files:Interrupt:Stop}
  \item \nameref{SEC:Files:Interrupt:Adjust_CFL}
  \item \nameref{SEC:Files:Interrupt:Start_TA}
  \item \nameref{SEC:Files:Interrupt:Dump_Restart}
  \item \nameref{SEC:Files:Interrupt:Dump_Solution}
 %\item \nameref{SEC:Files:Interrupt:Start_Reactions}
 %\item \nameref{SEC:Files:Interrupt:New_Restart}
\end{itemize}
%
with the description of each of the files provided in the following subsections.

%%%%%%%%%%%%%%%%%%%%%%%%%%%%%%%%%%%%%%%%%%%%%%%%%%%%%%%%%%%%%%%%%%%%%%%%%%%%%%%%

The value of $N$ depends on how GFR is compiled, and takes the following
values:
%
\\[1.5ex]
\begin{tabular}{ @{\qquad$N$ = } r @{\quad} p{0.8\linewidth} @{} }
  1 & when compiled in Debug mode \\
  50 & when compiled in Release\slash Production mode and \emph{disabling} code
  that accesses the PBS environment \\
  100 & when compiled in Release\slash Production mode and \emph{enabling} code
  that accesses the PBS environment
\end{tabular}
\\[1.5ex]

%%%%%%%%%%%%%%%%%%%%%%%%%%%%%%%%%%%%%%%%%%%%%%%%%%%%%%%%%%%%%%%%%%%%%%%%%%%%%%%%

These files are temporary in that they will automatically be deleted immediately
after each file is found, read, and processed.
%
This behavior is necessary because otherwise each file would continuously be
found and processed every $N$ time steps if not deleted.

%###############################################################################
%###############################################################################
%###############################################################################

\subsubsection{STOP\_GFR}
\label{SEC:Files:Interrupt:Stop}

This interrupt file allows the user to stop the currently running simulation,
either immediately or at a future iteration specified inside the file.
%
If this file is found, the first line of the file is read, and then the file is
closed and deleted.
%
If the value read from the first line is an integer, this value is saved as the
final iteration number for the running simulation.
%
Note that this value is treated as an \emph{aggregate} iteration number,
\thatis it directly corresponds to the first column of the iteration data
written to display\slash stdout.
%
This will cause the simulation to complete the standard shutdown sequence after
completing the final iteration, including writing the final restart and
solution output files.
%
If
  1) the \file{STOP\_GFR} file is empty;
  2) the read operation on the first line fails due to an invalid (\thatis
    non-integer) quantity; or
  3) the final iteration value read from the file is less than the current
    iteration number;
%
then GFR will immediately exit the solver loop and begin the standard shutdown
sequence.

%%%%%%%%%%%%%%%%%%%%%%%%%%%%%%%%%%%%%%%%%%%%%%%%%%%%%%%%%%%%%%%%%%%%%%%%%%%%%%%%
\vspace{1.0ex}
\noindent Using this interrupt file is as simple as running a command like
%
\\[1.2ex]\indent\texttt{echo 148234 > STOP\_GFR}\\[1.2ex]
%
%\begin{lstlisting}
%  echo 148234 > STOP_GFR
%\end{lstlisting}
%
on the command line in the simulation run directory.
%
This example would cause GFR to stop the simulation after completing the
148,234\supth{} time step.

%###############################################################################
%###############################################################################
%###############################################################################

\subsubsection{ADJUST\_CFL}
\label{SEC:Files:Interrupt:Adjust_CFL}

This interrupt file allows the user to modify the current CFL schedule of a
running simulation.
%
This file is assumed to contain the namelist \inptext{cfl\_adjustment} defined
as:
%
\begin{nmlblocklst}[
,label={LST:cfl_adjustment}
,caption={Namelist cfl\_adjustment.}
]
&cfl_adjustment
  CFL                    = [current simulation value]
  CFL_Ramping_Beg        = [current simulation value]
  CFL_Ramping_End        = [current simulation value]
  CFL_Ramping_Iterations = [current simulation value]
  CFL_Ramping_Start_Iter = [current simulation value]
/
\end{nmlblocklst}
%
The inputs in this namelist are identical to the respective inputs in
\sref{SEC:nml_timestep}.
%
Since this file is only read after a simulation has already started, each of
these five variables will have already been defined in the preprocessor when
the simulation was starting.
%
Hence, this is why the \quoterm{default value} for each input is given as%
\texttt{ = [current simulation value]}.

%%%%%%%%%%%%%%%%%%%%%%%%%%%%%%%%%%%%%%%%%%%%%%%%%%%%%%%%%%%%%%%%%%%%%%%%%%%%%%%%

If the \file{ADJUST\_GFR} file is found, GFR will attempt to read this namelist
from the file.
%
If any of these inputs are provided in the namelist, they will directly replace
the corresponding value within the simulation.
%
Once the current iteration number reaches the value specified in the input
\inp[nmlTimestep::CFLRampingStartIter]{CFL\_Ramping\_Start\_Iter},
the CFL ramping schedule is reset so that the CFL is linearly increased from
\inp[nmlTimestep::CFLRampingBeg]{CFL\_Ramping\_Beg} to
\inp[nmlTimestep::CFLRampingEnd]{CFL\_Ramping\_End} between iterations
$({\inptext{CFL\_Ramping\_Start\_Iter}}+1)$ and
$({\inptext{CFL\_Ramping\_Start\_Iter}}+%
\text{\inp[nmlTimestep::CFLRampingIterations]{CFL\_Ramping\_Iterations}})$.
%
If the value of {\inptext{CFL\_Ramping\_Start\_Iter}} read from the interrupt
file is less than the current iteration number, the CFL ramping schedule is
reset immediately.

%%%%%%%%%%%%%%%%%%%%%%%%%%%%%%%%%%%%%%%%%%%%%%%%%%%%%%%%%%%%%%%%%%%%%%%%%%%%%%%%

Whenever the CFL ramping schedule is changed by the inputs in this file, a new
restart file will be written just before the new schedule begins in case it is
too aggressive and causes the simulation to fail.

%%%%%%%%%%%%%%%%%%%%%%%%%%%%%%%%%%%%%%%%%%%%%%%%%%%%%%%%%%%%%%%%%%%%%%%%%%%%%%%%

This interrupt file can be very helpful to push the limits of finding the
largest CFL that can be used while also maintaining a stable solution.
%
One method to find this limit is to slowly keep pushing the CFL higher until the
simulation fails.
%
In past experience, plotting the data in \file{residual.dat} from the failed
simulation can identify an approximate iteration number where the solution began
to diverge.
%
This iteration number can be used with the iteration statistics written to the
display\slash stdout to identify the CFL value at that iteration.
%
From there, you should be able to run a stable simulation with a \quoterm{near
optimal} CFL by using a CFL that is 85\%--95\% of the CFL value when the
previous failed simulation began to diverge.

%###############################################################################
%###############################################################################
%###############################################################################

\subsubsection{START\_TIME\_AVERAGING}
\label{SEC:Files:Interrupt:Start_TA}

The simple presence of this interrupt file requests GFR to begin time averaging
the current simulation.
%
If the simulation is already time averaging, this interrupt file will be deleted
and ignored.
%
If the simulation is not time averaging, GFR will:
\begin{itemize}
  \item Begin collecting time averages of various variable quantities (see
    \sref{SEC:Files:Time_Averaging_Output}).
  \item Start writing time-averaged restart files
    \file{(Restart\_files\slash out.\#.ave)} each time an instantaneous restart
    file is written.
  \item Start writing time-averaged solution output files
    \file{(CGNS\_files\slash solution.time\_ave.\#.cgns)} each time an
    instantaneous solution output file is written.
  \item Reset the solution time written to the display\slash stdout to zero.
    From this point on, this quantity will report the total amount of time that
    the solution has been time averaged. Note, the actual simulation time is not
    reset within GFR and will continue accumulating in the background.
    Only the time quantity that is written to display is reset.
    %
    The iteration counter is also not reset.
  \item Reset the maximum L2 residual quantities used to determine
    convergence, and begin reporting the L2 residual of
    the \emph{time-averaged mean flow} variables to the \file{residual.dat}
    file instead of the \emph{instantaneous} variables.
\end{itemize}
%
This file requires no contents, and any contents will be ignored.
%
Using this interrupt file is as simple as running the following command on the
command line in the simulation run directory:
%
\\[1.2ex]\indent\texttt{touch START\_TIME\_AVERAGING}
%\begin{lstlisting}
%  touch START_TIME_AVERAGING
%\end{lstlisting}

%###############################################################################
%###############################################################################
%###############################################################################

\subsubsection{DUMP\_RESTART}
\label{SEC:Files:Interrupt:Dump_Restart}

The simple presence of this interrupt file requests GFR to output an
instantaneous restart file, as well as a time-averaged restart file if
currently time-averaging.
%
This file requires no contents, and any contents will be ignored.
%
Using this interrupt file is as simple as running the following command on the
command line in the simulation run directory:
%
\\[1.2ex]\indent\texttt{touch DUMP\_RESTART}
%\begin{lstlisting}
%  touch DUMP_RESTART
%\end{lstlisting}

%###############################################################################

\subsubsection{DUMP\_SOLUTION}
\label{SEC:Files:Interrupt:Dump_Solution}

The simple presence of this interrupt file requests GFR to output an
instantaneous solution output file, as well as a time-averaged output file if
currently time-averaging.
%
This file requires no contents, and any contents will be ignored.
%
Using this interrupt file is as simple as running the following command on the
command line in the simulation run directory:
%
\\[1.2ex]\indent\texttt{touch DUMP\_SOLUTION}
%\begin{lstlisting}
%  touch DUMP_SOLUTION
%\end{lstlisting}

%###############################################################################
%###############################################################################
%###############################################################################

%\subsubsection{START\_REACTIONS}
%\label{SEC:Files:Interrupt:Start_Reactions}

%###############################################################################
%###############################################################################
%###############################################################################

%\subsubsection{RESTART\_INTERRUPT}
%\label{SEC:Files:Interrupt:New_Restart}

%###############################################################################

